\scanpage{16}
\textbf{C. R. Acad. Sc. Paris, t. 278 (24 juin 1974)}\hfill\textbf{Série A — 1601}

\rule{\linewidth}{0.4pt}

$\mathrm T(\mathbf Z/2\mathbf Z,\mathrm R)$ définit dans (${}^5$) comme on le voit grâce au corollaire (7.4) de (${}^{11}$). D'après
le théorème 1 de (${}^5$), $\mathrm Q(\mathrm A)$ est un sous-groupe de $\mathrm Q(\mathrm K)$ car A est intégralement clos.
On conclut à l'aide du petit lemme des cinq\quad$\blacksquare$

\textsc{Lemme 2}. — \emph{L'application $a$ de $\operatorname{LW}(\mathrm A)$ dans $\operatorname{LW}(\mathrm K)$ est surjective.}

\emph{Preuve}. — D'après (${}^3$), théorème 5 et remarque, une forme quadratique Q sur K provient
d'une forme quadratique $q$ sur A si et seulement si les trois conditions suivantes sont
vérifiées :

(i) La dimension de Q est paire.

(ii) Son discriminant, i. e. le commutant de $\mathrm C_0(\mathrm Q)$ dans $\mathrm C(\mathrm Q)$ est de la forme $\mathrm S\underset{\mathrm A}{\otimes}\mathrm K$
où $\mathrm S\in\mathrm Q(\mathrm A)$.

(iii) Son algèbre de Clifford $\mathrm C(\mathrm Q)$ est de la forme $\mathrm D\underset{\mathrm A}{\otimes}\mathrm K$ où D est une A-algèbre
centrale séparable.

La condition (i) est vérifiée car la proposition 1 nous dit que si $\mathrm Q\in\operatorname{LW}(\mathrm K)$ elle est
le rang pair car ce groupe est isomorphe à $(8\mathbf Z)^r$.

Les conditions (ii) et (iii) sont immédiatement vérifiées car $[\mathrm C(\mathrm Q)]=[0]$.

Il existe donc une forme $q$ sur A telle que $b([q])=[\mathrm Q]$.

L'application $c$ étant injective, $[\mathrm C(q)]=[0]$ et donc $[q]\in\operatorname{LW}(\mathrm A)$\quad$\blacksquare$

\textsc{Proposition 2}. — \emph{L'application $a$ de $\operatorname{LW}(\mathrm A)$ dans $\operatorname{LW}(\mathrm K)$ est un isomorphisme et on
obtient la s. e. $0\to(8\mathbf Z)^r\to\operatorname{Witt}(\mathrm A)\to\operatorname{HW}(\mathrm A)\to0$ où $r$ est le nombre de valuations
réelles de} K.

\emph{Preuve}. — La proposition 2 de (${}^6$) montre que $b$ est injective. Il en est de même de $a$
et la proposition résulte des lemmes 1 et 2\quad$\blacksquare$

Nous avons la s. e. $0\to\mathrm N(\mathrm A)\to\operatorname{HW}(\mathrm A)\to\mathrm Q(\mathrm A)\to0$, où $\mathrm N(\mathrm A)$ est un sous-groupe
de $\operatorname{Br}^2(\mathrm A)$. A. Grothendieck a montré dans (${}^4$) que $\operatorname{Br}(\mathrm A)=(\mathbf Z/2\mathbf Z)^n$ où
$n=\operatorname{Max}\{0,r-1\}$. Soit $[\mathrm D]$ un élément de $\operatorname{Br}(\mathrm A)$. D'après le résultat rappelé au
paragraphe 2, l'algèbre $\mathrm D\underset{\mathrm A}{\otimes}\mathrm K$ est isomorphe à l'algèbre de Clifford $\mathrm C(\mathrm Q)$ d'une forme
de rang pair sur K. Les conditions (i), (ii), (iii) du théorème 5 de (${}^3$) sont donc vérifiées
et il existe une forme quadratique $q$ sur A telle que $\mathrm Q\cong q\underset{\mathrm A}{\otimes}\mathrm K$. Nous avons donc
\[
\mathrm D\underset{\mathrm A}{\otimes}\mathrm K\cong\mathrm C(\mathrm Q)\cong\mathrm C(q\underset{\mathrm A}{\otimes}\mathrm K)\cong\mathrm C(q)\underset{\mathrm A}{\otimes}\mathrm K.
\]

Comme $\operatorname{Br}(\mathrm A)$ s'injecte dans $\operatorname{Br}(\mathrm K)$, on a $[\mathrm D]=[\mathrm C(q)]$. Il en résulte que $\mathrm N(\mathrm A)=\operatorname{Br}(\mathrm A)$.
La suite exacte $0\to(\mathbf Z/2\mathbf Z)^n\to\operatorname{HW}(\mathrm A)\to\mathrm Q(\mathrm A)\to0$ ramène donc le calcul de $\operatorname{Witt}(\mathrm A)$
à celui de $\mathrm Q(\mathrm A)$.

\emph{Exemples} :

1\textsuperscript{o} $\mathrm A=\mathbf Z$ :
\[
\operatorname{Br}(\mathbf Z)=\mathrm Q(\mathbf Z)=0\qquad\text{et}\qquad r=1.
\]

D'où $\operatorname{HW}(\mathbf Z)=0$ et $\operatorname{Witt}(\mathbf Z)=8\mathbf Z$ [\emph{cf.} (${}^7$) et (${}^8$)].

\mspage{15}
